\abstractch

在分布式存储与云文件服务中，随着数据规模和并发量的增加，元数据处理链路往往会成为整个系统的瓶颈。传统元数据管理通常采用以 B+ 树为代表的页式索引结构，面对大量元数据操作请求时，会出现随机写、写放大和热点目录竞争等问题，限制吞吐增长。若将元数据存储切换为基于 LSM-tree（Log-Structured Merge-tree）的键值存储，并据此重构关键文件操作流程，更新粒度可细化到单条记录，随机写也可以转化为顺序批量写，是缓解元数据瓶颈的一条可行路径。

基于此方向构建了文件系统 RucksFS，元数据采用键值存储，并实现了核心文件操作。系统由元数据服务器、数据服务器和客户端三部分组成，将元数据管理与文件数据管理分离。元数据服务器负责管理 inode 属性、目录项索引等核心内容，后端采用 RocksDB 存储，按 inodes、dir\_entries 等列族组织数据；并发控制采用悲观并发控制事务，保证文件操作的原子性和一致性。针对共享父目录热点，还引入了 DeltaOp 增量更新与懒折叠机制，减少父目录 inode 的就地改写和锁争用。数据服务器负责文件数据管理与落盘，提供读写、截断、删除和刷盘接口。客户端实现文件操作和跨服务编排，基于 FUSE（Filesystem in Userspace）承接内核文件操作请求，向上提供 POSIX（Portable Operating System Interface）标准文件语义，并将请求路由到服务器进行处理。整体上，系统将元数据键值存储、并发控制与核心文件操作整合为一体，形成了可验证的工程实现。

正确性方面，使用 pjdfstest 进行验证，实现范围内的 182 个可执行用例全部通过。性能方面，使用 mdtest 在 1 台 64 核服务器和多台客户端组成的集群上进行测试。在 32 并发下，RucksFS 的 create 吞吐相对 NFS v4.2（ext4 后端）提升约 4.6 倍。在更高并发下验证 DeltaOp 增量更新与懒折叠机制：在 384 并发下，create 与 remove 吞吐分别达到未启用版本的 2.96 倍和 2.63 倍。这些结果说明，采用键值存储并结合 LSM-tree 与 DeltaOp 的方案，在保证语义正确性的同时，能够有效提升高并发元数据操作的可扩展性。

\vspace{5pt}

\keywordsch 文件系统；元数据管理；键值存储；日志结构合并树；增量更新；悲观并发控制；用户态文件系统